% !TeX spellcheck = en_US
% !TeX encoding = UTF-8
% !TeX TXS-program:compile = txs:///latexmk/[-pdf -silent -shell-escape -latexoption="-synctex=1"]

% Workaround for bug, see https://tex.stackexchange.com/questions/575325/options-not-working-if-custom-class-is-in-subfolder
\makeatletter\def\input@path{{./docstyle/}}\makeatother

% IDSCreport in article mode: generic short document without chapters,
% table of contents, or information page
\documentclass[article,notoc]{IDSCreport}

% The class prints the nomenclature with \chapter*, which does not exist in
% article mode. A proposal needs no nomenclature, so skip it.
\renewcommand{\typesetnomenclature}{}

% Acronyms: \ac prints "long form (SHORT)" at first use, then only "SHORT"
\newacronym{auv}{AUV}{autonomous underwater vehicle}
\newacronym{rmse}{RMSE}{root-mean-square error}
\newacronym{rtk}{RTK}{real-time kinematic}

\title{Dynamics- and Tether-Aware Mission Planner for the POLARIS AUV}
\type{Bachelor's Thesis Proposal}
\author{Noel Bühler}
\date{November 18, 2026}

\begin{document}

\section{Context and Problem}
\label{sec:ContextAndProblem}

In Project POLARIS, ARIS Space, IDSC, and D-EAPS at ETH Zurich developed a
tethered \ac{auv} that measures ice thickness from below. The \ac{auv} is
deployed through a single hole in the ice and stays connected to the surface by
a fiber-optic tether up to \qty{300}{\meter} long. The mission objective is to
cover a \qty{10000}{\meter\squared} area with an upward-facing imaging SONAR,
take point measurements while touching the ice over a subset of that area, and
return safely to the hole.

The current autonomy stack plans straight lines between waypoints and tracks
them with a modified Pure Pursuit controller that switches into a dedicated
turning mode at corners. With this stack, the \ac{auv} completed a
\qtyproduct{20 x 20}{\meter} grid at Lake Zurich in \qty{28}{\minute} with a
cross-track \ac{rmse} of \qty{0.45}{\meter}.

The current planner, however, computes the path from geometry alone and ignores
how the path affects the vehicle and the mission. This matters in three areas.
The first is vehicle dynamics. The \ac{auv} accelerates, brakes, and turns
slowly relative to its cruise speed, so these maneuvers take up a large share
of the mission time. In the field tests, this caused corner overshoot, missed
goal tolerances, and replanning at \qty{90}{\degree} turns.
The second is the tether and the battery. Every turn adds twist to the tether,
and every meter traveled consumes energy that must stay in reserve for the
return to the hole.
The third is localization. Under the ice, the \ac{auv} navigates by dead
reckoning, and its position and heading estimates drift over time. The vehicle
can only correct this drift by rising to the ice and obtaining a GNSS fix
through it with the dual-antenna \ac{rtk} receiver. Between fixes, the heading
is integrated from the gyroscope alone and drifts more during maneuvers, so yaw
action directly degrades localization. To keep the drift bounded, the \ac{auv}
must obtain a fix approximately every \qty{300}{\meter} traveled.
% Source: POLARIS final report, localization appendix: open-trajectory drift
% 0.42 m per 100 m (grid_02), L_max = 384 m at the 2 m bound (one-sided 95%),
% 300 m adopted as the operational dead-reckoning budget.

\paragraph{Problem Statement} The \ac{auv} needs a mission planner that turns a
set of waypoints or a survey area into a path the vehicle can actually follow,
and that schedules GNSS fixes often enough to keep the dead-reckoning drift
within the mission requirement of \qty{\pm 2}{\meter}. The planner should keep
mission time, energy use, tether twist, and heading drift low, while always
leaving enough battery to return to the hole.

\paragraph{Research Question} How can an offline mission planner generate
dynamically feasible, tether- and drift-aware, velocity-stamped paths for
underwater waypoint missions, and how much does it improve tracking accuracy,
mission time, energy consumption, tether twist, and localization compared to
straight-line planning?

\section{Approach}
\label{sec:Approach}

The thesis develops a 2D, mission-level planner that computes offline a
dynamically feasible, velocity-stamped path through given waypoints inside a
specified safe area. The safe area is assumed to be free of obstacles. The
planner is based on a kinematic vehicle model, i.e.\ a unicycle with velocity,
yaw-rate, and acceleration limits, and on an empirical model of the electrical
current drawn for a given acceleration. All limits are parameters. They are
initialized from the Lake Zurich data and updated as the hardware and the
path- and velocity-tracking controller developed in the partner thesis by Ridh
evolve, so the planner stays valid without redesign.

Two waypoint maneuvers are distinguished:
\begin{description}
    \item[Stop waypoints:] the vehicle comes to rest and may turn on the
          spot, so the segments between them are straight lines.
    \item[Continuous waypoints:] the vehicle passes without stopping, so
          the path must respect forward motion with bounded yaw rate and yaw
          acceleration.
\end{description}

The planner minimizes mission time, energy use, and two distinct yaw
measures. \emph{Net yaw} sets the tether twist. Turns in opposite directions
cancel, so the goal is zero net yaw. \emph{Total absolute yaw} drives the
gyroscope drift and never cancels: a path that turns \qty{+360}{\degree} and
then \qty{-360}{\degree} has no tether twist but a lot of accumulated heading motion. The
planner is constrained by the safe area, the velocity and acceleration limits
of the vehicle, the maximum electrical current, and the deployed tether length.
It is built in stages, each of which is usable on its own:
\begin{enumerate}
    \item \textbf{Waypoint generator:} continuous waypoints for boustrophedon
          SONAR coverage of the survey area, and stop waypoints for the point
          measurements.
    \item \textbf{Waypoint ordering:} a traveling-salesman-type ordering of the
          waypoints that trades off path length, net yaw, and total yaw, with
          nearest neighbor and 2-opt as reference heuristics.
    \item \textbf{Corner smoothing:} curvature-continuous transitions at
          continuous waypoints that respect the yaw-rate and yaw-acceleration
          limits and stay inside the safe area.
    \item \textbf{Velocity profile:} a speed for every point of the path under
          the acceleration and current limits, trading mission time against
          energy. The controller receives a velocity-stamped path, and the
          planner integrates the profile only internally to estimate the
          mission time.
    \item \textbf{Return planner and battery watchdog:} a path back to the hole
          that avoids the area already covered by the tether, and a continuous
          check that this return still fits within the remaining battery
          capacity.
    \item \textbf{Fix scheduling:} using the same path integration with an
          assumed drift rate, the planner predicts the localization drift along
          the path and inserts ice-touch fix maneuvers as stop waypoints before
          the localization budget is exceeded. The existing behavior tree
          executes the maneuver, while the planner schedules it and accounts for
          its time and energy.
    \item \textbf{Online fallback (optional):} a lightweight version of the
          planner for replanning during the mission. It replans after large
          deviations in position or velocity, or when the localization
          uncertainty tracked by the state estimator calls for an earlier fix
          maneuver. In addition, it compares the measured remaining battery
          capacity with the energy estimated for a safe return to the hole along
          a path that avoids the area covered by the tether, and it triggers the
          return before this margin is used up.
\end{enumerate}

\paragraph{Evaluation} The main comparison keeps the controller fixed and
swaps only the planner, so differences can be attributed to the planner, as
listed in \cref{tab:Configurations}. Pure Pursuit cannot follow a velocity
profile, so pairing it with the new planner would understate the benefit of
the planner. If this tracking controller is not ready in time, a simple
path-tracking controller with velocity feedforward replaces it in simulation.

\begin{table}[htbp]
    \centering
    \caption{Planner and Controller Configurations Compared in the Evaluation}
    \label{tab:Configurations}
    \begin{tabular}{ll}
        \toprule
        Configuration                        & Role \\
        \midrule
        Straight-line planner + Pure Pursuit & Reference: current validated stack \\
        Straight-line planner + tracking controller & Baseline: isolates the controller \\
        New planner + tracking controller           & Proposed system \\
        \bottomrule
    \end{tabular}
\end{table}

The metrics are grouped by the aspects the planner aims to improve:
\begin{description}
    \item[Tracking accuracy:] cross-track \ac{rmse} and corner overshoot.
    \item[Mission time:] total mission time.
    \item[Energy:] total energy drawn from the battery, i.e.\ battery life, and
          peak current.
    \item[Tether:] net yaw, i.e.\ tether twist, and maximum distance to the hole.
    \item[Localization:] total absolute yaw and number of fix maneuvers.
\end{description}
The configurations are compared in simulation and, if possible, in field tests
on the vehicle. The test missions are defined during the thesis.

\section{Scope and Plan}
\label{sec:ScopeAndPlan}

\paragraph{In Scope} Offline 2D planning at constant depth in an obstacle-free
safe area, stages~1 to~6 above, integration into the existing ROS~2 and
\texttt{nav2} stack, and validation in the Gazebo and ArduSub simulation against
the current stack. If time runs short, fix scheduling (stage~6) is the first
item to move to the extensions.

\paragraph{Out of Scope} Obstacle avoidance, tether entanglement around
obstacles, 3D planning, and the design of the path-tracking controller
itself.

\paragraph{Optional Extensions} The online fallback (stage~7), including
uncertainty-triggered fix maneuvers based on the covariance of the state
estimator, field tests on the vehicle, a tether map that estimates where the
tether lies and feeds the return planner, and SONAR lane spacing derived from
the operating depth and the beam angle of the imaging SONAR.

\paragraph{Plan} \Cref{tab:Plan} shows the planned schedule.

\begin{table}[htbp]
    \centering
    \caption{Planned Schedule}
    \label{tab:Plan}
    \begin{tabular}{lp{0.75\textwidth}}
        \toprule
        Weeks  & Work package \\
        \midrule
        1--2   & Literature review, waypoint generator and waypoint ordering (stages~1 and~2) \\
        3--4   & System identification: kinematic limits and current model from field tests \\
        5--6   & Corner smoothing (stage~3) \\
        7--8   & Velocity profile with acceleration and current limits (stage~4) \\
        9      & Return planner and battery watchdog (stage~5) \\
        10     & Fix scheduling (stage~6) \\
        11--12 & Simulation evaluation against the reference and the baseline, and a field test if possible \\
        13--14 & Thesis writing and buffer \\
        \bottomrule
    \end{tabular}
\end{table}

\end{document}
